\section{Block exchange on a guessed active set}
\label{sec:exchange}

Given a candidate set $\Active$, Proposition~\ref{prop:sub} gives the working-set
solution directly:
\begin{equation}
  \bigl(C_\Active\T G^{-1} C_\Active\bigr)\lambda_\Active
   = b_\Active - C_\Active\T x_u,
  \qquad
  x = x_u + G^{-1}C_\Active \lambda_\Active,
  \qquad x_u = G^{-1}a,
  \label{eq:pdassolve}
\end{equation}
with $\lambda_i = 0$ for $i \notin \Active$. Both solves go through one Cholesky
factorisation of $G$, computed once, and one of the $k \times k$ dual Hessian,
computed per repair. The repair rule reads the two conditions that
\eqref{eq:pdassolve} leaves open, \eqref{eq:dfeas} and \eqref{eq:pfeas}, as
instructions:
\begin{equation}
  \Active' = \{1,\dots,m_{\mathrm{eq}}\} \;\cup\;
  \bigl\{\, i > m_{\mathrm{eq}} \;:\;
    (i \in \Active \text{ and } \lambda_i \ge 0)
    \ \text{ or }\
    (c_i\T x < b_i) \,\bigr\},
  \label{eq:repair}
\end{equation}
that is: an active constraint whose multiplier has gone negative does not belong
in the set and leaves it; an inactive constraint that is violated at $x$ does
belong and joins it; equalities are always in. If $\Active' = \Active$ then
\eqref{eq:pfeas}--\eqref{eq:comp} hold to the tolerance used, and
\eqref{eq:pdassolve} supplies \eqref{eq:stat} by construction, so the pair is
optimal by Proposition~\ref{prop:sufficient}.

The starting guess is the one the rule itself dictates when nothing has been tried:

\begin{equation}
  \Active_0 = \{1,\dots,m_{\mathrm{eq}}\} \;\cup\;
  \bigl\{\, i > m_{\mathrm{eq}} \;:\; c_i\T x_u < b_i \,\bigr\},
  \label{eq:start}
\end{equation}
the equalities together with whatever the unconstrained minimiser violates, which
is already the answer when it violates nothing.

Two features of this iteration matter for what follows. The whole set is exchanged
at once, so the number of repairs is not tied to the number of constraints that
have to change status, and in practice it is nearly independent of $n$
(Section~\ref{sec:measurement}). And the set is named rather than reached: nothing
in \eqref{eq:repair} or \eqref{eq:start} consults the geometry of the normals it
selects.

\begin{remark}[The contrast with the walk]\label{rem:walk}
A dual active-set method modifies $\Active$ by one index per step, and the step
rules of~\cite{goldfarb1983} are exactly what make $C_\Active$ full column rank
for free. A constraint whose normal lies in the span of those already held admits
no primal movement, so its step is limited by the multiplier ratio test rather
than by the primal one, and it can be added only after enough partial steps have
removed the dependence. Linear dependence is therefore unreachable, and no rank
test appears anywhere in such an implementation. This is the protection that
guessing forfeits, and it is worth naming as a property rather than an accident,
because it explains why the difficulty of Section~\ref{sec:obstruction} has no
counterpart in the classical method.
\end{remark}
